\section*{Hoofdstuk \ref{ch:b2:plant-plasticity} --- Aanpassingen van planten en fenotypische plasticiteit}

\begin{solution}{exo:b2:plant-plasticity:1}
Plasticiteit: één \omterm{def:b2:meiosis-heredity:vocabulary}{genotype}, meerdere \omterm{def:b2:meiosis-heredity:vocabulary}{fenotypen} naargelang de omgeving
(zonne- en schaduwbladeren aan één eik). \omterm{def:b2:plant-plasticity:plasticity}{Reactienorm}: de kromme van het
\omterm{def:b2:meiosis-heredity:vocabulary}{fenotype} van een \omterm{def:b2:meiosis-heredity:vocabulary}{genotype} tegen een omgevingsvariabele (bladdikte tegen
licht). \omterm{def:b2:plant-plasticity:plasticity}{Acclimatisatie}: een omkeerbare fysiologische aanpassing (de
winterharding van rogge in de herfst). \omterm{def:b2:plant-plasticity:plasticity}{Adaptatie}: een erfelijke afstemming,
over generaties geselecteerd (het CAM-metabolisme van een cactus).
\end{solution}

\begin{solution}{exo:b2:plant-plasticity:2}
Intacte coleoptiel die van één kant wordt belicht: buigt naar het licht. Top
afgedekt met een ondoorzichtig kapje: geen buiging --- de top neemt het
licht waar. Top afgedekt met een doorzichtig kapje: buigt --- het kapje zelf
is onschadelijk. Basis afgeschermd door een ondoorzichtig kokertje, top
bloot: buigt --- de waarneming zit in de top en de respons eronder.
\end{solution}

\begin{solution}{exo:b2:plant-plasticity:3}
\omterm{prop:b2:plant-plasticity:sunshade}{Zonneblad}: klein, dik, twee of drie palissadelagen, dikke cuticula, veel
huidmondjes, hoge maximale fotosynthese en een hoog compensatiepunt.
\omterm{prop:b2:plant-plasticity:sunshade}{Schaduwblad}: groot, dun, één palissadelaag, meer chlorofyl per gram, lage
ademhaling, laag compensatiepunt. De keuze wordt gemaakt terwijl het blad
een primordium in de knop is, op grond van het licht dat het ontvangt; een
volgroeid blad kan niet meer veranderen.
\end{solution}

\begin{solution}{exo:b2:plant-plasticity:4}
De sluitcellen nemen kalium en water op, zwellen op en buigen, omdat hun
binnenwanden dikker zijn, uit elkaar, waardoor de porie opengaat. Blauw
licht en een laag inwendig $\mathrm{CO_2}$ openen ze 's ochtends;
\omterm{def:b2:plant-flowering:dormancy}{abscisinezuur}, uit wortels in uitdrogende grond of uit het blad, sluit ze
bij droogte, net als zeer droge lucht.
\end{solution}

\begin{solution}{exo:b2:plant-plasticity:5}
$\varphi = R/(R + 0.77\,FR)$ met $FR = 1$: $1.15$: $0.60$; $0.7$:
$0.48$; $0.2$: $0.21$; $0.05$: $0.06$. De drempel $0.4$ valt tussen de
open plek ($0.48$) en het bladerdak ($0.21$): de \omterm{thm:b2:plant-plasticity:phytochrome}{schaduwvermijding} schakelt
aan onder het bladerdak.
\end{solution}

\begin{solution}{exo:b2:plant-plasticity:6}
$E = 0.25\times 0.025 = \qty{6.25}{mmol.m^{-2}.s^{-1}}$; $A =
0.25\times 100/1.6 = \qty{15.6}{\micro mol.m^{-2}.s^{-1}}$; $A/E =
2.5\times 10^{-3}$ mol koolstof per mol water, d.w.z.\ 400 watermoleculen
per koolstof, $400\times 18/12 = \qty{600}{g}$ water per gram koolstof. In
12 uur: $6.25\times 10^{-3}\times 43200 =
\qty{270}{mol/m^2}$, \qty{4.9}{L} per vierkante meter.
\end{solution}

\begin{solution}{exo:b2:plant-plasticity:7}
$E = 0.25\times 0.005 = \qty{1.25}{mmol.m^{-2}.s^{-1}}$, $A$
onveranderd: $A/E = 0.0125$, 80 watermoleculen per koolstof,
\qty{120}{g/g} --- een vijfde van de kosten overdag.
\end{solution}

\begin{solution}{exo:b2:plant-plasticity:8}
$v = 2r^{2}\Delta\rho g/9\eta = 2\times 4\times 10^{-12}\times
400\times 9.81/0.09 = \qty{3.5e-7}{m/s}$: \qty{0.35}{\micro m/s}, ongeveer
\qty{30}{s} om de cel over te steken. Op een klinostaat verandert de
richting van de zwaartekracht ten opzichte van de cel om de paar seconden
van het zakken, zodat de korrels zich nooit aan één kant ophopen en de
gemiddelde prikkel over de integratietijd van de plant (minuten) nul is.
\end{solution}

\begin{solution}{exo:b2:plant-plasticity:9}
Gemiddelde verlenging in twee uur $0.04$; kanten $1.2\times 0.04 = 0.048$
en $0.8\times 0.04 = 0.032$; $\theta = 20\times 0.016/2 = 0.16$
rad, ongeveer \qty{9}{\degree}. De buiging neemt toe met $0.08$ rad per uur;
$\pi/2$ vergt ongeveer \qty{20}{h}.
\end{solution}

\begin{solution}{exo:b2:plant-plasticity:10}
De reserves volstaan voor \qty{50}{cm} stengel. Brandnetelveld:
\qty{30}{cm} in 7.5 dagen voor \qty{30}{mg} --- de zaailing bereikt het licht
met reserves over. Bos: \qty{10}{m} ligt buiten bereik; de zaailing besteedt
zijn \qty{50}{mg} in twaalf dagen aan een halve meter bleke stengel en sterft
geëtioleerd, terwijl een schaduwtolerante strategie --- enkele dunne
bladeren, trage groei --- hem jarenlang in leven had kunnen houden.
\end{solution}

\begin{solution}{exo:b2:plant-plasticity:11}
IJs vormt zich eerst in de extracellulaire ruimten, waar de oplossing
verdunder is; de dampspanning boven ijs is lager dan boven het water van de
cel, zodat water de cel verlaat om zich bij het ijs te voegen en de cel
uitdroogt en haar membranen ineenstorten. Een week kou laadt de cel met
suikers en beschermende eiwitten die haar vriespunt verlagen en membranen
stabiliseren, en laat antivrieseiwitten maken die de kristalgroei beperken.
Het programma het hele jaar aan houden, leidt koolstof af van groei naar
bescherming en vertraagt de deling: een blijvend geharde plant is een kleine
plant.
\end{solution}

\begin{solution}{exo:b2:plant-plasticity:12}
Het verroodsignaal is eerlijk wanneer de buur kan worden overgroeid en een
leugen wanneer het een boom is; de warme week is eerlijk in april en een
leugen in februari; in een dicht gewas zijn de buren van elke plant haar
gelijken, zodat het signaal eerlijk is in vorm maar nutteloos, en veredelaars
selecteren planten die het negeren. Een \omterm{def:b2:plant-plasticity:plasticity}{reactienorm} is het register van hoe
vaak elk signaal in het verleden van de lijn de waarheid sprak, en hij faalt
telkens wanneer het heden van dat verleden verschilt.
\end{solution}

\begin{solution}{pb:b2:plant-plasticity:1}
\textbf{1.} Open veld: $1.15/(1.15 + 0.77) = 0.60$; bladerdak: $0.2/(0.2 +
0.77) = 0.21$.
\textbf{2.} Lage $\varphi$: planten erboven --- strekken. Onder een
schaduwdoek blijft $\varphi$ op $0.60$: een donkere dag, geen buur; geen
\omterm{thm:b2:plant-plasticity:phytochrome}{schaduwvermijding}, alleen tragere groei door gebrek aan licht.
\textbf{3.} $1 + 2(0.60 - 0.21) = 1.78$: \qty{2.7}{cm} per dag.
\textbf{4.} \qty{37}{cm} tegenover \qty{21}{cm}.
\textbf{5.} Buigen naar de open plek (fototropie, via de blauwlichtreceptor
fototropine) en snellere strekking en heroriëntatie van de bladeren aan de
kant met de hogere $R{:}FR$ (via fytochroom).
\textbf{6.} Schaduwbladeren worden groot en dun gebouwd om schaars licht
goedkoop op te vangen; in de volle zon raken ze oververhit, ondergaan ze
fotoinhibitie en verbranden ze, en moet de plant ze door zonnebladeren
vervangen.
\textbf{7.} \qty{2.7}{cm} per dag over een zone van \qty{15}{mm} is
\qty{0.11}{cm} per uur, een relatieve verlenging van $0.074$ per uur;
beschaduwde kant $1.3\times 0.074 = 0.096$, belichte kant $0.7\times 0.074 =
0.052$.
\textbf{8.} $\theta = 15\times(0.096 - 0.052)/2 = 0.33$ rad, ongeveer
\qty{19}{\degree} in een uur.
\textbf{9.} \qty{60}{\degree} is $1.05$ rad: ongeveer drie uur.
\textbf{10.} $v = 2\times 6.25\times 10^{-12}\times 400\times
9.81/0.09 = \qty{5.5e-7}{m/s}$, \qty{0.55}{\micro m/s}: \qty{22}{s}
om de cel over te steken.
\textbf{11.} De fototropie, en de \omterm{thm:b2:plant-plasticity:phytochrome}{schaduwvermijding} vergroot zelf de hoek
waaronder de scheut tevreden groeit: licht is de hulpbron waarom wordt
gevochten, en een stengel die zich tegen de zwaartekracht in rechtte, zou
terug de schaduw in draaien.
\textbf{12.} De wortel buigt omlaag tot hij verticaal staat: statolieten
zakken naar de onderkant, auxine hoopt zich daar op, en in een wortel remt
de extra auxine de groei, zodat de bovenkant sneller groeit. Dezelfde
herverdeling, een tegengestelde gevoeligheid van de cellen.
\textbf{13.} $E = 0.2\times 0.015 = \qty{3}{mmol.m^{-2}.s^{-1}}$; over
$2\times 10^{-3}$ \unit{m^2} en \qty{50400}{s}: \qty{0.30}{mol},
\qty{5.4}{g} water per dag.
\textbf{14.} $A = 0.2\times 80/1.6 = \qty{10}{\micro mol.m^{-2}.s^{-1}}$;
per dag $0.50$ mol koolstof per vierkante meter; $A/E = 3.3\times
10^{-3}$: 300 watermoleculen per koolstof, \qty{450}{g/g}.
\textbf{15.} $0.50\times 2\times 10^{-3} = \qty{1}{mmol}$: \qty{12}{mg}
koolstof per dag.
\textbf{16.} Open plek: $E = \qty{6}{mmol.m^{-2}.s^{-1}}$, \qty{10.9}{g} per
dag; efficiëntie \qty{900}{g/g}.
\textbf{17.} $E = 0.02\times 0.03 = \qty{0.6}{mmol.m^{-2}.s^{-1}}$
(tienvoudig lager); $A = 0.02\times 200/1.6 = \qty{2.5}{\micro
mol.m^{-2}.s^{-1}}$ (viervoudig lager). Het water daalt het meest: naarmate
het blad zijn inwendige $\mathrm{CO_2}$ verlaagt, wordt de gradiënt voor
koolstof groter, terwijl die voor water dat niet kan.
\textbf{18.} \qty{4.25}{g} water, waarvan \qty{0.85}{g} mag worden
verloren; bij \qty{10.9}{g} per dag ($0.78$ g per uur) verwelkt hij in
ongeveer een uur.
\textbf{19.} CAM is traag en vraagt succulent opslagweefsel; een zaailing die
onder een bladerdak naar het licht racet, waar het waterverlies toch klein
is, heeft snelheid nodig, geen zuinigheid.
\textbf{20.} $100/2 = \qty{50}{cm}$.
\textbf{21.} Brandnetelveld: \qty{40}{cm} in 15 dagen voor \qty{80}{mg} ---
hij bereikt het licht. Bos: \qty{8}{m} is onbereikbaar; hij besteedt zijn
reserves aan een halve meter en sterft. Beter in het bos: kort blijven, dunne
schaduwbladeren maken en op een open plek wachten --- de manier van de beuk.
\textbf{22.} Berk: een steile \omterm{def:b2:plant-plasticity:plasticity}{reactienorm}, met sterke strekking bij lage
$\varphi$, omdat de buren van een pionier even groot zijn als hijzelf. Beuk:
een vlakke \omterm{def:b2:plant-plasticity:plasticity}{reactienorm}, weinig strekking en veel tolerantie, omdat zijn buren
bomen zijn.
\textbf{23.} Schakel de schaduwvermijdingsrespons uit --- houd fytochroom B
actief of verwijder de \omterm{prop:b2:cell-differentiation:transcriptional}{transcriptiefactoren} die het in toom houdt --- en
controleer bloeitijdstip, kieming en stengelsterkte, die door dezelfde route
worden gestuurd.
\textbf{24.} Een plant die haar buren niet kan waarnemen, wordt overgroeid
telkens wanneer ze had kunnen winnen; een plant die het signaal altijd
gelooft, verspilt zichzelf onder elke boom. De winnende \omterm{def:b2:plant-plasticity:plasticity}{reactienorm} zet op
het signaal in verhouding tot hoe vaak het in het verleden van de lijn gelijk
had.
\textbf{25.} $\varphi = 0.60$ in het open veld en $0.21$ onder het bladerdak;
\qty{37}{cm} na 14 dagen; \qty{5.4}{g} water per dag in de schaduw,
\qty{10.9}{g} op de open plek.
\end{solution}
